\chapter{The Requisites Mode}\index{modes!Requisites Mode|textbf}
\label{ch:requisites}

\begin{versebox}
Those states that generate another state ---\\
those are its requisites, its ground.\\
A seed, a fire, a lamp ---\\
without the cause, no fruit is found.
\end{versebox}

\noindent Nothing arises from nothing. That is the fundamental axiom of Buddhist thought, and the Requisites Mode is the tool that maps it.\footnote{Pali: \textit{``Tattha katamo parikkhāro hāro? `Ye dhammā yaṃ dhammaṃ janayantī\,''ti.''} The mnemonic fragment is ``those states that generate that state'' (\textit{ye dhammā yaṃ dhammaṃ janayanti}). The word \textit{parikkhāra} means ``requisite, accessory, equipment'' --- what you need in order for something to arise. Whatever state generates another state, the first is the \textit{parikkhāra} (requisite) of the second. Its defining characteristic: the characteristic of generating (\textit{janakalakkhaṇa}).}

Whatever generates another thing is that thing's requisite. But there is a crucial distinction to be made. Two things generate: a \textit{cause} and a \textit{condition}.\index{cause and condition|textbf} And they are not the same.

A cause is specific --- it belongs uniquely to its effect. A condition is shared --- it supports many effects indiscriminately.\footnote{Pali: \textit{``Dve dhammā janayanti hetu ca paccayo ca. Tattha kiṃlakkhaṇo hetu, kiṃlakkhaṇo paccayo? Asādhāraṇalakkhaṇo hetu, sādhāraṇalakkhaṇo paccayo.''} Cause (\textit{hetu}) = specific, unshared (\textit{asādhāraṇa}). Condition (\textit{paccaya}) = common, shared (\textit{sādhāraṇa}). This distinction is fundamental to the Netti's causal analysis.}

\section*{The Seed and the Soil}

Think of a seed sprouting into a plant. The seed is the cause --- specific, irreplaceable. Only a rice seed produces rice; only a mango seed produces a mango tree. But the earth and water are conditions --- shared, generic. The same earth and water support rice, mango, and every other plant. The seed is what makes the effect \textit{this particular one}. The earth and water are what make any growth possible at all.\footnote{Pali: \textit{``Yathā kiṃ bhave? Yathā aṅkurassa nibbattiyā bījaṃ asādhāraṇaṃ, pathavī āpo ca sādhāraṇā. Aṅkurassa hi pathavī āpo ca paccayo sabhāvo hetu.''} The Netti's first analogy: for the arising of a sprout (\textit{aṅkura}), the seed (\textit{bīja}) is the specific cause; earth (\textit{pathavī}) and water (\textit{āpo}) are the shared conditions.}

Or think of a lamp. The dish, the wick, and the oil are conditions --- they support the flame. But they cannot produce light on their own. You cannot light a dish. The fire is the cause --- the specific, irreplaceable element that makes light happen. The dish, wick, and oil are necessary. But they are not sufficient.\footnote{Pali: \textit{``Yathā vā pana thālakañca vaṭṭi ca telañca padīpassa paccayabhūtaṃ na sabhāvahetu, na hi sakkā thālakañca vaṭṭiñca telañca anaggikaṃ dīpetuṃ padīpassa paccayabhūtaṃ. Padīpo viya sabhāvo hetu hoti.''} The lamp analogy: dish (\textit{thālaka}), wick (\textit{vaṭṭi}), and oil (\textit{tela}) = conditions. Fire (\textit{aggi}) = the essential cause (\textit{sabhāvo hetu}). The Netti then draws four parallel distinctions: essential nature = cause, external nature = condition; internal = cause, external = condition; generating = cause, supporting = condition; specific = cause, shared = condition.}

So: the essential nature is the cause; the external nature is the condition. What is internal is the cause; what is external is the condition. What generates is the cause; what supports is the condition.

\section*{Milk and Curd}

The Netti adds a startling observation. When milk is placed in a pot and becomes curd, you never see milk and curd existing at the same time. The milk \textit{becomes} the curd. Similarly, the cause and the condition do not coexist with their effect. The cause transforms into the effect; the condition supports the transformation. But you never catch all three in the same frame.\footnote{Pali: \textit{``Yathā vā pana ghaṭe duddhaṃ pakkhittaṃ dadhi bhavati, na catthi ekakālasamavadhānaṃ duddhassa ca dadhissa ca. Evamevaṃ natthi ekakālasamavadhānaṃ hetussa ca paccayassa ca.''} The milk-curd analogy. ``There is no simultaneous co-occurrence (\textit{ekakālasamavadhāna}) of milk and curd. Likewise, there is no simultaneous co-occurrence of cause and condition [with their effect].'' This is a statement about temporal succession in causation: causes produce effects by transforming into them, not by standing alongside them.}

This is not abstract philosophy. This is the engine of continued existence. The cycle of rebirth arose with cause and condition. The Buddha said: dependent on ignorance, formations arise\index{formations}; dependent on formations, consciousness arises --- and so on through the entire twelve-link chain of dependent origination. Every link is both the cause of the next and the condition for the ones that follow.\footnote{Pali: \textit{``Ayañhi saṃsāro sahetu sappaccayo nibbatto. Vuttaṃ hi avijjāpaccayā saṅkhārā, saṅkhārapaccayā viññāṇaṃ, evaṃ sabbo paṭiccasamuppādo.''} The cycle of existence (\textit{saṃsāra}) arose with cause and condition. The twelve links of dependent origination are the prime example of the Requisites Mode in action.}

\section*{The Two Kinds of Cause}

Even within ``cause,'' the Netti distinguishes two types: the immediate cause and the indirect cause. Previous ignorance --- the underlying tendency, lying dormant in the mind --- is the immediate cause of the ignorance-obsession that flares up in active experience. Like a seed producing a sprout --- directly, intimately, from one moment to the next.\footnote{Pali: \textit{``Purimikā avijjā pacchimikāya avijjāya hetu. Tattha purimikā avijjā avijjānusayo pacchimikā avijjā avijjāpariyuṭṭhānaṃ, purimiko avijjānusayo pacchimikassa avijjāpariyuṭṭhānassa hetubhūto paribrūhanāya, bījaṅkuro viya samanantarahetutāya.''} Two types of cause: immediate cause (\textit{samanantarahetu}) = the underlying tendency (\textit{anusaya}) of ignorance, which directly produces the active obsession (\textit{pariyuṭṭhāna}) of ignorance, ``like a seed producing a sprout.'' Indirect cause (\textit{paramparahetu}) = the chain of effects that follows remotely.}

But there is also the indirect cause --- the chain of consequences that follows at a distance. The tree that grew from the seed eventually produces fruit. The seed did not directly produce the fruit; it produced the sprout, which produced the trunk, which produced the branches, which produced the fruit. The seed is the indirect cause of the fruit. Both types of causation are at work in the cycle of existence.

\section*{The Chain of Implications}

The Netti then maps a remarkable chain --- a sequence of logical implications that shows how a single break in continuity would unravel the entire cycle of suffering:

Where there is non-interruption, there is continuity. Where there is continuity, there is arising. Where there is arising, there is fruit. Where there is fruit, there is reconnection. Where there is reconnection, there is renewed existence. Where there is renewed existence, there is obstruction. Where there is obstruction, there is obsession. Where there is obsession, there is non-uprooting. Where there is non-uprooting, there is underlying tendency. Where there is underlying tendency, there is non-penetration. Where there is non-penetration, there is ignorance. Where there is ignorance, there is consciousness --- uncomprehended, functioning as a seed.\footnote{Pali: \textit{``Avupacchedattho santati attho, nibbatti attho phalattho, paṭisandhi attho punabbhavattho, palibodhattho pariyuṭṭhānattho, asamugghātattho anusayattho, asampaṭivedhattho avijjattho, apariññātattho viññāṇassa bījattho. Yattha avupacchedo tattha santati \ldots yattha avijjā tattha sāsavaṃ viññāṇaṃ apariññātaṃ, yattha sāsavaṃ viññāṇaṃ apariññātaṃ tattha bījattho.''} Twelve causal implications: non-interruption → continuity → arising → fruit → reconnection → renewed existence → obstruction → obsession → non-uprooting → underlying tendency → non-penetration → ignorance → consciousness-as-seed. Each link is the requisite of the next. The chain shows how the cycle is self-sustaining: ignorance produces consciousness-as-seed, which produces non-interruption, which starts the chain again. Break any link, and the chain collapses.}

Twelve implications, each the requisite of the next. The chain is circular: ignorance produces consciousness-as-seed, which produces non-interruption, which starts the chain again. Break any link, and the whole thing collapses.

\section*{The Upward Chain}

But the Requisites Mode does not only map bondage. It also maps liberation. The five training-aggregates form their own causal chain: virtue is the requisite of concentration\index{virtue}\index{concentration}. Concentration is the requisite of wisdom. Wisdom is the requisite of liberation. Liberation is the requisite of the knowledge-and-vision of liberation.\footnote{Pali: \textit{``Sīlakkhandho samādhikkhandhassa paccayo, samādhikkhandho paññākkhandhassa paccayo, paññākkhandho vimuttikkhandhassa paccayo, vimuttikkhandho vimuttiñāṇadassanakkhandhassa paccayo.''} Five training-aggregates as a causal chain: virtue → concentration → wisdom → liberation → knowledge-and-vision of liberation. Each is the requisite of the next. This is the path mapped not as a list of factors but as a sequence of causes.}

And another, more subtle chain: knowing the assembly leads to knowing rapture. Knowing rapture leads to knowing what has been attained. Knowing what has been attained leads to knowing oneself.\footnote{Pali: \textit{``Titthaññutā pītaññutāya paccayo, pītaññutā pattaññutāya paccayo, pattaññutā attaññutāya paccayo.''} A chain of self-knowledge: knowing the assembly (\textit{titthaññutā}) → knowing rapture (\textit{pītaññutā}) → knowing what's been attained (\textit{pattaññutā}) → knowing oneself (\textit{attaññutā}). These are among the seven ``knowings'' (\textit{sattaññā}) of a good monk.}

\section*{Every Link Is Double}

The Netti closes with a sweeping demonstration. Take the simplest sensory event: dependent on the eye and visible forms, eye-consciousness arises. The eye is the condition in terms of predominance. The visible forms are the condition in terms of object. Light is the condition in terms of support. But attention --- attention is the cause proper. Without attention, the eye sees nothing. With attention, seeing happens.\footnote{Pali: \textit{``Yathā vā pana cakkhuñca paṭicca rūpe ca uppajjati cakkhuviññāṇaṃ. Tattha cakkhu ādhipateyyapaccayatāya paccayo, rūpā ārammaṇapaccayatāya paccayo. Āloko sannissayatāya paccayo, manasikāro sabhāvo hetu.''} The fourfold causal analysis of sense-cognition: eye = predominance-condition (\textit{ādhipateyyapaccayatā}), forms = object-condition (\textit{ārammaṇapaccayatā}), light = support-condition (\textit{sannissayatā}), attention (\textit{manasikāra}) = essential cause (\textit{sabhāvo hetu}).}

And then the Netti runs through the entire twelve-link chain of dependent origination one more time, this time labeling each link as both condition and essential cause: formations are the condition of consciousness \textit{and} its essential cause. Consciousness is the condition of name-and-form \textit{and} its essential cause. And so on, all the way through sorrow, lamentation, pain, grief, and despair.\footnote{Pali: \textit{``Saṅkhārā viññāṇassa paccayo, sabhāvo hetu. Viññāṇaṃ nāmarūpassa paccayo, sabhāvo hetu \ldots Domanassaṃ upāyāsassa paccayo, sabhāvo hetu. Evaṃ yo koci upanissayo sabbo so parikkhāro.''} Every link in the chain of dependent origination serves as both condition and essential cause of the next. The closing formula: ``whatever support there is, all of that is a requisite'' (\textit{yo koci upanissayo sabbo so parikkhāro}).}

Whatever support there is --- all of it is a requisite. That is the Requisites Mode: the map of what generates what, the chart of causal dependencies that shows you exactly where to intervene if you want the chain to stop.\footnote{Pali: \textit{``Tenāha āyasmā mahākaccāyano `ye dhammā yaṃ dhammaṃ janayantī\,''ti.''}}
